% ID: SEL-D0019-OB-E02
% ESTADO: validada. S020, cotejada con examen y criterios oficiales.
\textbf{a)} En los puntos del dominio, la regla del cociente da
$f'(x)=a/(bx+1)^2$. El paralelismo impone $f'(0)=2$, luego $a=2$.
Por otra parte, $f(-1)=1$ exige $b\ne1$ y $-2/(1-b)=1$.
Resulta $\boxed{a=2,\ b=3}$, con $-1$ en el dominio.

\textbf{b)} Con $a=b=1$ se obtiene $f(x)=1-1/(x+1)$.
El límite por la izquierda de $-1$ es $+\infty$ y por la derecha es
$-\infty$, de modo que $\boxed{x=-1}$ es una asíntota vertical.
Al tender $x$ a cualquiera de los dos infinitos, $1/(x+1)$ tiende a
cero. Por tanto $\boxed{y=1}$ es la asíntota horizontal.
